%% ma: L10-B14-0014
%% lop: 10
%% chuong: 5
%% bai: 14
%% dang: 10.14.4
%% loai: TL
%% muc: TH
%% dap_an: "-"
%% nguon: "chuyen-de-cac-so-dac-trung-cua-mau-so-lieu-khong-ghep-nhom-toan-10-chuong-trinh-moi.pdf (GV. Phan Nhật Linh) (P04), Bài 03 Dạng 1 Phần I Câu 4"
%% kien_thuc: "Lớp 10 Bài 14 (khoảng biến thiên)"
%% trang_thai: chinh-thuc
%% ghi_chu: "Câu nằm ở Phần I nhưng đề gốc không có phương án, đã nhập dạng tự luận. Tự duyệt theo yêu cầu giáo viên 22/10/2026."
\begin{ex}
Điểm số của hai vận động viên bắn cung trong $10$ lần bắn được ghi lại như sau:
\begin{center}
\begin{tabular}{|c|c|c|c|c|c|c|c|c|c|c|}\hline
Vận động viên A & $10$ & $9$ & $8$ & $10$ & $9$ & $9$ & $9$ & $10$ & $9$ & $8$\\\hline
Vận động viên B & $5$ & $10$ & $10$ & $10$ & $10$ & $7$ & $9$ & $10$ & $10$ & $10$\\\hline
\end{tabular}
\end{center}
Tính khoảng biến thiên của mỗi dãy số liệu. Căn cứ trên chỉ số này, vận động viên nào có thành tích bắn ổn định hơn?
\loigiai{$R_A=10-8=2$; $R_B=10-5=5$.

Vì $R_A<R_B$ nên theo khoảng biến thiên, vận động viên A có thành tích ổn định hơn.}
\end{ex}
